%=============================================================================
% EXPERIMENT 4 APPENDIX SECTION -- copy and paste this whole block.
%
% Self-contained: no custom macros, no \input, nothing to define elsewhere.
%
% Requires in your preamble:
%   \usepackage{graphicx}
%   \usepackage{placeins}   % for \FloatBarrier
%
% Figure paths read "figures/experiment-4/<...>", relative to the main
% document, matching the convention used for experiment 2. Copy
% outputs/finals/experiment-4/{probe,representation}/ into
% <maindoc>/figures/experiment-4/ (or point \graphicspath there). To relocate,
% find-and-replace the single prefix "figures/experiment-4/" throughout.
%
% Source: frozen experiment-4 snapshot, assembled 2026-09-08 by
% experiments/finalize_experiment_4.py at commit f5b1212c.
% See README.md and MANIFEST.csv.
%=============================================================================

% --- Float placement. Defaults restored at the end of this block, so pasting
% --- it mid-document does not affect other sections.
\renewcommand{\topfraction}{0.92}
\renewcommand{\bottomfraction}{0.7}
\renewcommand{\textfraction}{0.06}
\renewcommand{\floatpagefraction}{0.6}

%=============================================================================
\section{Experiment 4: encoder representation}
\label{app:experiment-4-figures}
%=============================================================================

This appendix reports the diagnostics for the \texttt{vsvfixed} variant of the
attentional drift-diffusion model, in which both the base drift $v$ and its
dispersion $s_v$ are pinned, leaving five estimated cells: the drift advantage
$v_{\mathrm{diff}}$, the attentional weight $\omega$, the threshold $a$, the
non-decision time $\tau$, and the non-decision dispersion $s_\tau$.

Two arms are compared throughout, trained and evaluated on shared validation
data so that they differ only in how the observations are summarized. The
\emph{frozen} arm pairs a frozen CogFormer encoder with a FlowMatching
posterior; the \emph{baseline} arm replaces that encoder with a SetTransformer
summary learned for this task, keeping the same posterior network. The question
is whether a summary transferred from the model-class CogFormer is as good as,
and encodes the same thing as, one learned for the task directly.

Figures~\ref{fig:e4-vsvfixed-baseline-ecdf} and~\ref{fig:e4-vsvfixed-frozen-ecdf}
give the posterior-level check: rank-ECDF calibration for each arm.
Figure~\ref{fig:e4-vsvfixed-rdm-rsa} then gives the representation-level check,
comparing the two arms' representational geometry directly.

\FloatBarrier

%-----------------------------------------------------------------------------
\subsection{Posterior calibration}
\label{sec:e4-vsvfixed-ecdf}
%-----------------------------------------------------------------------------

\begin{figure}[htbp]
  \centering
  \includegraphics[width=\linewidth]{figures/experiment-4/probe/vsvfixed/baseline/calibration_ecdf.pdf}
  \caption[vsvfixed, baseline arm: ECDF calibration]{\textbf{Attentional DDM,
    \texttt{vsvfixed} target, baseline (learned SetTransformer summary) ---
    ECDF calibration.} Rank-ECDF difference (empirical minus uniform) with
    simultaneous 95\% confidence bands, one panel per estimated parameter. A
    calibrated posterior stays inside the band; excursions above or below
    indicate bias, and a systematically shallow or steep curve indicates over-
    or under-dispersion. Calibration here should be read alongside the
    contraction values: $\omega$ is barely identified in this design
    (contraction $0.068$, NRMSE $0.972$), so its posterior is close to the
    prior and calibrates almost trivially, whereas $a$ and $\tau$ are sharply
    contracted ($0.990$ and $0.977$) and their bands are the informative test.}
  \label{fig:e4-vsvfixed-baseline-ecdf}
\end{figure}

\begin{figure}[htbp]
  \centering
  \includegraphics[width=\linewidth]{figures/experiment-4/probe/vsvfixed/frozen/calibration_ecdf.pdf}
  \caption[vsvfixed, frozen arm: ECDF calibration]{\textbf{Attentional DDM,
    \texttt{vsvfixed} target, frozen (transferred CogFormer encoder) --- ECDF
    calibration.} As Figure~\ref{fig:e4-vsvfixed-baseline-ecdf}, for the arm
    whose summary statistics come from the frozen CogFormer encoder rather than
    a task-specific learned summary. Both arms are calibrated, and the pattern
    across parameters is the same; calibration alone therefore does not
    separate them. The paired comparison on shared validation data does: a
    classifier two-sample test scores $0.623$ against an empirical null of
    $0.501$, an excess of $+0.122$ at $z = 29.7$, with the learned summary
    ahead on log-density and CRPS.}
  \label{fig:e4-vsvfixed-frozen-ecdf}
\end{figure}

\FloatBarrier
\clearpage

%-----------------------------------------------------------------------------
\subsection{Representational geometry}
\label{sec:e4-vsvfixed-rsa}
%-----------------------------------------------------------------------------

\begin{figure}[htbp]
  \centering
  \includegraphics[width=\linewidth]{figures/experiment-4/representation/intercept_only_vsvfixed/rdm_rsa.pdf}
  \caption[vsvfixed: RDM representational similarity analysis]{\textbf{
    Attentional DDM, \texttt{vsvfixed} target --- representational similarity
    analysis.} Representational dissimilarity matrices (RDMs) over $200$
    parameter draws, computed with the crossnobis distance from $4$ independent
    repeats per draw. \emph{Left to right:} the RDM of the learned
    SetTransformer summary, the RDM of the frozen CogFormer encoder, their
    residual, and a scatter of the two off-diagonal vectors against each other.
    Both matrices are z-scored over their off-diagonal entries, because
    crossnobis units depend on the representation and the two arms are roughly
    an order of magnitude apart in raw scale; z-scoring is the transform
    Pearson correlation already applies, so the maps show exactly what the
    reported $r$ measures. One consequence is that the dissimilarities are
    signed and centred on the mean rather than on zero, so the colour scale
    reads ``more or less dissimilar than typical'' and the raw ``$0$ =
    identical'' anchor is not preserved. The observed correspondence is
    $r = 0.791$, against a dimensionality-matched null of $0.840 \pm 0.052$
    obtained by random $128 \rightarrow 16$ projection --- that is, at or below
    what an uninformative projection achieves. The two arms are not encoding
    the same geometry, so the posterior-level gap is not merely a difference of
    decoder.}
  \label{fig:e4-vsvfixed-rdm-rsa}
\end{figure}

\FloatBarrier

% --- Restore LaTeX default float parameters -----------------------------------
\renewcommand{\topfraction}{0.7}
\renewcommand{\bottomfraction}{0.3}
\renewcommand{\textfraction}{0.2}
\renewcommand{\floatpagefraction}{0.5}

%=============================================================================
% END OF EXPERIMENT 4 APPENDIX SECTION
%=============================================================================
